%%%  Bakalářská práce - Roth: Kompilátor jazyka Forth v Rustu
%%%  Albert Klinkovský, 2026


%%  Pro získání PDF souboru dokumentu je třeba tento zdrojový text v
%%  LaTeXu přeložit (dvakrát) programem pdfLaTeX.

%%  V případě použití programu BibLaTeX pro tvorbu seznamu literatury
%%  je poté ještě třeba spustit program Biber s parametrem jméno
%%  souboru zdrojového textu bez přípony a následně opět (dvakrát)
%%  přeložit zdrojový text programem pdfLaTeX.

\documentclass[
%  master,
%  program=ainfvs,
%  printversion,
  biblatex,
%  language=english,
%  font=sans,
  figures=false,
%  tables=false,
%  theorems,
  sourcecodes,
  glossaries,
%  index
]{kidiplom}

%% Název práce
\title{Návrh a implementace zásobníkového jazyka s JIT kompilací v jazyce Rust}
\title[english]{Design and implementation of a stack language with JIT compilation in Rust}

%% Autor a vedoucí
\author{Albert Klinkovský}
\supervisor{Mgr. Jan Laštovička, Ph.D.}

%% Anotace
\annotation{Práce se zabývá návrhem a implementací kompilátoru
  zásobníkového programovacího jazyka inspirovaného jazykem Forth.
  Implementace v jazyce Rust podporuje generování kódu do jazyků
  Rust a C. Součástí práce je interaktivní prostředí REPL založené
  na průběžné kompilaci a dynamickém načítání knihoven, optimalizační
  průchody nad mezikódem a standardní knihovna jazyka.}

\annotation[english]{This thesis deals with the design and implementation
  of a compiler for a stack-based programming language inspired by Forth.
  The Rust implementation supports code generation to Rust and C.
  The work includes an interactive REPL environment based on incremental
  compilation and dynamic library loading, optimization passes over
  intermediate representation, and a standard library for the language.}

%% Klíčová slova
\keywords{kompilátor; Forth; Rust; zásobníkový jazyk; JIT kompilace; optimalizace}
\keywords[english]{compiler; Forth; Rust; stack-based language; JIT compilation; optimization}

%% Poděkování
\thanks{Děkuji vedoucímu práce za odborné vedení a cenné rady.}

%% Bibliografie
\bibliography{bibliografie.bib}

%% Další balíky
\usepackage{tikz}
\usetikzlibrary{shapes,arrows,positioning}
\usepackage{booktabs}

%% listings does not ship Rust/Forth/bash definitions by default.
%% Define minimal languages used in code examples.
\lstdefinelanguage{Rust}{
  sensitive=true,
  morecomment=[l]{//},
  morecomment=[s]{/*}{*/},
  morestring=[b]",
  keywords={
    as,break,const,continue,crate,else,enum,extern,false,fn,for,if,impl,in,
    let,loop,match,mod,move,mut,pub,ref,return,self,Self,static,struct,
    super,trait,true,type,unsafe,use,where,while
  }
}

\lstdefinelanguage{Forth}{
  sensitive=true,
  morecomment=[s]{(}{)},
  morestring=[b]"
}

\lstdefinelanguage{bash}{
  sensitive=true,
  morecomment=[l]{\#},
  morestring=[b]"
}

%% kidiplom.cls may call \printglossary even when glossaries are disabled.
%% Provide no-op fallbacks to keep builds non-interactive.
\providecommand{\printglossary}[1][]{ }
\providecommand{\printglossaries}{ }

\begin{document}
\maketitle

%% Pomocná makra
\newcommand{\Roth}{\textsc{Roth}}
\newcommand{\Forth}{\textsc{Forth}}

\input{chapters/01-uvod}


%% ===================================================================
%% KAPITOLA 2: TEORETICKÝ ZÁKLAD
%% ===================================================================
\section{Teoretický základ}
\label{sec:teorie}

\input{chapters/02-forth}

\input{chapters/02-kompilatory}

\input{chapters/02-rust}


%% ===================================================================
%% KAPITOLA 3: ANALÝZA A NÁVRH
%% ===================================================================
\section{Analýza a návrh}
\label{sec:navrh}

\input{chapters/03-navrh}


%% ===================================================================
%% KAPITOLA 4: IMPLEMENTACE
%% ===================================================================
%% ===================================================================
%% KAPITOLA 4: IMPLEMENTACE
%% ===================================================================
\section{Implementace}
\label{sec:implementace}

\subsection{Lexikální analýza}
\label{sec:lexer}

Lexikální analyzátor převádí vstupní text na sekvenci tokenů, které jsou
následně zpracovány parserem. V~kompilátoru \\Roth{} je lexer implementován
v~modulu \texttt{src/lexer.rs} jako struktura \texttt{Lexer}, která si uchovává:
(1) celý vstup jako řetězec, (2) aktuální pozici ve vstupu a~(3) informace o~řádku
a~sloupci pro diagnostiku chyb.

Každý token nese kromě svého typu také původní text (pole \texttt{raw}) a~pozici
\texttt{Position} (řádek, sloupec a~offset ve vstupu). Tato volba zjednodušuje
generování chybových hlášení ve všech následujících fázích.

Typy tokenů jsou definovány výčtem \texttt{TokenType} v~\texttt{src/types.rs}.
Lexer rozpoznává celočíselné literály (\texttt{i32}), identifikátory (slova),
oddělovače definic (\texttt{:} a~\texttt{;}), komentáře v~závorkách
\texttt{(\dots)} a~řetězcové literály.

Specifickým prvkem inspirovaným \Forth{} je podpora syntaxe \texttt{S\" ...\"},
která umožňuje zapisovat řetězec bez explicitního uzavírání do tokenů jednotlivých
slov. Lexer tento vzor detekuje ještě před obecným rozlišením tokenu.

\begin{lstlisting}[language=Rust, caption={Rozpoznání tokenu v lexeru}, float=htbp]
fn next_token(&mut self, start_pos: Position) -> Result<Token, ParseError> {
    let ch = self.current_char();
    if (ch == 'S' || ch == 's') && self.peek_char() == Some('"') {
        return self.read_s_quote_literal(start_pos);
    }
    match ch {
        '(' => self.read_comment(start_pos),
        '"' => self.read_string_literal(start_pos),
        ':' | ';' => self.read_definition_token(start_pos),
        _ => self.read_token(start_pos),
    }
}
\end{lstlisting}

Metoda \texttt{read\_token} načte nejdelší podřetězec oddělený bílými znaky
a~pokusí se jej převést na číslo. Pokud převod selže, token je interpretován jako
slovo a~normalizován na velká písmena (\texttt{to\_uppercase}). To sjednocuje
zpracování slov nezávisle na velikosti písmen ve zdrojovém kódu.

Komentáře v~závorkách jsou tokenizovány jako \texttt{TokenType::Comment}.
Parser i~další fáze je následně mohou ignorovat nebo využít (např. pro zachování
informativních poznámek v~IR ve formě \texttt{Comment} instrukcí).

\subsection{Syntaktická analýza}
\label{sec:parser}

Parser (\texttt{src/parser.rs}) je implementován jako jednoduchý
rekurzivně-sestupný analyzátor nad vektorem tokenů. Vzhledem
k~tomu, že zásobníkový jazyk má lineární zápis (program je sekvence slov
a~literálů), je výsledný abstraktní syntaktický strom (AST) záměrně mělký.

Kořenem AST je uzel \texttt{Program}, který obsahuje seznam příkazů. Jedinou
hierarchickou konstrukcí jsou definice slov uzavřené mezi \texttt{:} a~\texttt{;},
které jsou reprezentovány uzlem \texttt{Definition} s~vlastním tělem.

Základní uzly AST jsou definovány v~\texttt{src/types.rs}.

\begin{lstlisting}[language=Rust, caption={Základní uzly AST}]
pub enum AstNode {
    Number(i32, Position),
    Word(String, Position),
    StringLiteral(String, Position),
    Definition { name: String, body: Vec<AstNode>, position: Position },
    VariableDeclaration { name: String, position: Position },
    Program(Vec<AstNode>),
}
\end{lstlisting}

Parser v~celém programu ignoruje tokeny typu \texttt{Comment}. Tím se snižuje
komplexita syntaktické analýzy bez ztráty informace o~pozicích chyb, protože
pozice je nesena i~ostatními tokeny.

Definice slova je parsována následovně: po \texttt{:} musí následovat jméno
slova, poté parser postupně načítá uzly těla až do \texttt{;}. Vnořené definice
jsou zakázány a~jsou vyhodnoceny jako syntaktická chyba.

Parser zároveň implementuje speciální pravidlo pro konstrukci \texttt{VARIABLE}.
Pokud se v~toku tokenů objeví slovo \texttt{VARIABLE}, parser očekává následující
identifikátor, který uloží do uzlu \texttt{VariableDeclaration}. Vlastní práce
s~proměnnými je řešena až v~nižších fázích (sémantika a~IR).

\subsection{Sémantická analýza}
\label{sec:analyzer}

Sémantická analýza (\texttt{src/analyzer.rs}) ověřuje, že program používá pouze
existující slova a~že uživatelské definice neporušují pravidla jazyka.
V~kompilátoru \\Roth{} je analýza založena na udržování tří tabulek:
\texttt{builtin\_words}, \texttt{defined\_words} a~\texttt{defined\_variables}.

Při inicializaci analyzátoru je tabulka vestavěných slov naplněna seznamem
primitivních operací (aritmetika, práce se zásobníkem, I/O, řízení toku,
práce s~pamětí). Vestavěná slova jsou považována za součást definice jazyka a~nelze
je redefinovat.

Analýza prochází AST rekurzivně:
\begin{itemize}
	\item při uzlu \texttt{Definition} ověří, že jméno není vestavěné, následně
	      přidá slovo mezi definovaná ještě \emph{před} analýzou těla;
	      tím je umožněna rekurze pomocí \texttt{RECURSE},
	\item při uzlu \texttt{Word} ověří, že se jedná o~vestavěné slovo, uživatelsky
	      definované slovo, nebo proměnnou,
	\item při \texttt{VariableDeclaration} registruje proměnnou v tabulce.
\end{itemize}

V~režimu REPL je tato fáze rozšířena o~kontext předchozích vstupů: před voláním
\texttt{analyze} jsou do analyzátoru doplněny již definované slova a~proměnné,
aby bylo možné na ně navazovat (viz sekce~\ref{sec:repl}).

\subsection{Generování mezikódu}
\label{sec:ir}

Mezikód (IR) je generován v~modulu \texttt{src/ir\_lowering.rs} transformací
AST na strukturu \texttt{IRProgram}. Implementace používá pomocnou strukturu
\texttt{IRBuilder} (\texttt{src/ir.rs}), která umožňuje postupně emitovat
instrukce do právě budované funkce.

Lowering je realizován ve třech krocích, které odpovídají přirozené struktuře
\Forth{}-like programu:
\begin{enumerate}
	\item nejprve se sesbírají všechny definice slov do mapy \texttt{word\_definitions},
	\item následně se pro každé slovo vytvoří samostatná IR funkce,
	\item nakonec je přeložena hlavní část programu (kód mimo definice) do funkce \texttt{main}.
\end{enumerate}

Číselné literály jsou překládány na instrukci \texttt{Push(Constant)}.
Řetězcové literály jsou v~současné implementaci reprezentovány jako sekvence
znaků uložených na zásobník (ASCII kódy) následovaná délkou řetězce; to
odpovídá očekávanému stack efektu slova \texttt{TYPE}.

\paragraph{Proměnné.}
Při zpracování uzlu \texttt{VariableDeclaration} je proměnné přidělena celočíselná
adresa z~rostoucího čítače. Při použití proměnné jako slova (tj. token \texttt{Word}
odpovídá jménu proměnné) lowering vloží na zásobník její adresu.
Operace \texttt{@} a~\texttt{!} jsou mapovány na instrukce \texttt{Load} a~\texttt{Store}.

\paragraph{Řídicí struktury.}
Základní struktury \texttt{IF/ELSE/THEN} a~\texttt{DO/LOOP} jsou překládány na
instrukce se skoky a~návěstími (\texttt{Jump}, \texttt{JumpIfNot}, \texttt{Label})
pomocí generování unikátních návěstí v~\texttt{IRBuilder}. Pro správné párování
konstrukcí lowering udržuje pomocné zásobníky (\texttt{loop\_stack} a~\texttt{conditional\_stack}).

\paragraph{Výpočet stack effect.}
Po vygenerování IR je nad funkcemi spuštěna analýza stack efektu
(\texttt{StackEffectAnalyzer}), která z~anotací instrukcí (procedurální makro
\texttt{\#[derive(StackEffect)]} v~\texttt{roth-derive}) vypočítá čistou změnu
hloubky zásobníku pro každou funkci. Tato informace je později využita
optimalizátorem i~pro ladicí výpisy.

\subsection{Optimalizace}
\label{sec:optimalizace}

\subsubsection{Constant folding}

Průchod constant folding (\texttt{ConstantFoldingPass} v~\texttt{src/ir\_optimizer.rs})
vyhledává lokální vzory typu \texttt{Push(a) Push(b) op} a~nahrazuje je jednou
instrukcí \texttt{LoadConst(result)}. Implementace podporuje aritmetické,
relační i~logické operace. Pro booleovské výsledky je použita konvence \Forth{}:
\texttt{-1} znamená pravdu a~\texttt{0} nepravdu.

\subsubsection{Dead code elimination}

Eliminace mrtvého kódu (\texttt{DeadCodeEliminationPass}) odstraňuje instrukce,
které nemají vliv na výsledek výpočtu. V~aktuální implementaci se jedná zejména
o~\texttt{Nop} a~krátké sekvence jako \texttt{Push \; Drop} nebo \texttt{Dup \; Drop},
které vznikají jako vedlejší produkt jiných optimalizací.

\subsubsection{Peephole optimalizace}

Peephole optimalizace (\texttt{PeepholeOptimizationPass}) pracuje nad
malými okny instrukcí (typicky délky 2--3). Cílem je nahradit časté
sekvence efektivnějším ekvivalentem, případně vytvořit příležitost
pro následné odstranění instrukcí (např. nahrazení dvojice instrukcí
\texttt{Dup Drop} za \texttt{Nop}, který odstraní DCE).

\subsubsection{Strength reduction}

Strength reduction (\texttt{StrengthReductionPass}) implementuje standardní
algebraické identity a~zjednodušení, například násobení dvěma na \texttt{Dup Add}
nebo eliminaci neutrálních prvků (\texttt{+ 0}, \texttt{* 1}).

\subsubsection{Inlining}

Inlining (\texttt{FunctionInliningPass}) nahrazuje instrukci \texttt{Call} tělem
volané funkce, pokud je funkce dostatečně malá (výchozí limit 20 instrukcí)
a~neobsahuje skoky ani rekurzivní volání. Průchod je implementován jako
globální optimalizace nad celým \texttt{IRProgram}, protože potřebuje přístup
k~mapě všech funkcí.

Optimalizační pipeline (\texttt{IROptimizer}) spouští průchody iterativně, dokud
nedojde ke stabilizaci, nebo dokud není překročen maximální počet iterací
(výchozí hodnota 10). Iterativní spouštění je důležité například proto, že
inlining může vytvořit nové konstantní výrazy vhodné pro constant folding.

\subsection{Generování cílového kódu}
\label{sec:codegen}

\subsubsection{Backend pro Rust}

Generování cílového kódu pro Rust je implementováno dvěma cestami.
První cesta (\texttt{rust-ir}) používá přímočarý generátor
\texttt{IRRustGenerator} (\texttt{src/ir\_codegen.rs}), který převádí IR do
samostatného spustitelného programu v~Rust. Vygenerovaný kód obsahuje strukturu
\texttt{OptimizedForth} s~datovým zásobníkem \texttt{Vec<i32>} a~metodami,
které implementují jednotlivá uživatelská slova jako vnořené funkce.

Pro instrukce řídicího toku generátor používá jednoduchý stavový automat
nad čítačem instrukcí (\texttt{\_\_pc}), což umožňuje reprezentovat skoky
\texttt{Jump}/\texttt{JumpIf}/\texttt{JumpIfNot} i~bez explicitních \texttt{goto}.

Druhá cesta je novější modulární framework (\texttt{src/codegen/*}), který
odděluje překlad instrukcí (\texttt{IRTranslator}) od emitování kódu
(\texttt{CodeEmitter}) a~od specifik cílového jazyka (\texttt{TargetLanguage}).
V~této práci je tento framework využit především pro experimenty a~rozšiřitelnost.

\subsubsection{Backend pro C}

Backend pro jazyk C (\texttt{c-ir}) je implementován generátorem
\texttt{IRCGenerator} v~\texttt{src/ir\_codegen.rs}. Vygenerovaný kód
implementuje zásobník jako pole pevné velikosti \texttt{STACK\_SIZE} a~používá
pomocné funkce \texttt{push}/\texttt{pop}. Tato varianta je vhodná pro jednodušší
portaci a~snadné spuštění na systémech bez toolchainu pro Rust.

Stejně jako u~Rust backendu je možné použít i~modulární variantu backendu
(\texttt{c-modular}), která je registrována v~\texttt{codegen/registry.rs}.

\subsection{REPL s JIT kompilací}
\label{sec:repl}

REPL je implementován v~\texttt{src/repl/} a~realizuje tzv. \emph{JIT-like}
přístup: každý vstup uživatele je přeložen do Rust kódu, zkompilován do
sdílené knihovny (\texttt{cdylib}) a~načten za běhu programu.

Vstupní řetězec prochází stejnými fázemi jako kompilace souboru:
lexing, parsing, sémantická analýza, IR lowering a~optimalizace. Následně se
použije specializovaný generátor \texttt{ReplCodegen} (\texttt{src/repl/codegen.rs}),
který vytvoří kód exportující funkci \texttt{\_\_repl\_entry}. Ta přijímá
\texttt{RuntimeContext} z~knihovny \texttt{roth-runtime} a~vykoná instrukce nad
perzistentním zásobníkem a~pamětí.

Kompilace a~dynamické načítání jsou zapouzdřeny v~\texttt{LibraryLoader}
(\texttt{src/repl/loader.rs}). Loader zapisuje vygenerovaný Rust do dočasného
adresáře, spustí \texttt{rustc --crate-type=cdylib} a~výslednou knihovnu načte
pomocí knihovny \texttt{libloading}. REPL si uchovává všechny načtené knihovny
v~paměti, aby nedošlo k~uvolnění symbolů.

Perzistentní stav mezi příkazy je rozdělen na:
\begin{itemize}
	\item \textbf{runtime kontext} (\texttt{RuntimeContext}) --- zásobník,
	      paměť proměnných a~registry slov,
	\item \textbf{kompilační kontext} (\texttt{CompilerContext}) --- slovník již
	      definovaných slov v~IR a~množinu deklarovaných proměnných. Tyto
	      informace umožňují sémantické ověření a~IR lowering s~ohledem na
	      historii REPL.
\end{itemize}

\subsection{Standardní knihovna}
\label{sec:stdlib}

Standardní knihovna je uložena ve složce \texttt{std/} a~je psaná přímo
v~jazyce \\Roth{}. Hlavní soubor \texttt{std/std.rt} postupně načítá jednotlivé
moduly pomocí direktivy \texttt{INCLUDE}.

Knihovna je rozdělena do tematických souborů:
\begin{itemize}
	\item \texttt{core.rt} --- základní konstanty a~utility,
	\item \texttt{stack.rt} --- rozšířené zásobníkové operace,
	\item \texttt{math.rt} --- matematické funkce,
	\item \texttt{io.rt} --- formátovaný výstup a~I/O pomocné funkce,
	\item \texttt{control.rt} --- rozšířené řídicí konstrukce,
	\item \texttt{compiler.rt} --- dokumentace slov implementovaných na úrovni kompilátoru.
\end{itemize}

V~offline kompilaci je direktiva \texttt{INCLUDE} zpracována ještě před lexikální
analýzou (funkce \texttt{preprocess\_includes} v~\texttt{src/main.rs}). Preprocesor
rekurzivně rozbalí všechny inkluze, detekuje cykly a~vloží do výsledného zdroje
značky ve formě komentářů, aby bylo možné zpětně dohledat původ vloženého kódu.




%% ===================================================================
%% KAPITOLA 5: TESTOVÁNÍ A VYHODNOCENÍ
%% ===================================================================
%% ===================================================================
%% KAPITOLA 5: TESTOVÁNÍ A VYHODNOCENÍ
%% ===================================================================
\section{Testování a vyhodnocení}
\label{sec:testovani}

\subsection{Testovací metodika}

Testování kompilátoru bylo rozděleno do dvou úrovní:
\begin{itemize}
	\item \textbf{jednotkové testy} ověřují jednotlivé moduly (lexer, parser,
	      sémantická analýza, IR, optimalizace a~codegen) na malých izolovaných
	      případech,
	\item \textbf{integrační testy} ověřují celý překladový řetězec přes
	      veřejné CLI rozhraní a~kontrolují, že je vygenerovaný kód vytvořen
	      a~obsahuje očekávané konstrukce.
\end{itemize}

Jednotkové testy jsou implementovány ve dvou formách:
(1) testy přímo v modulech v~\texttt{src/} pomocí \texttt{\#[cfg(test)]} a
(2) testy v~adresáři \texttt{tests/}, kde jsou jednotlivé části kompilátoru
testovány přes veřejné API knihovny \texttt{roth}.

Integrační testy (\texttt{tests/integration\_tests.rs}) spouštějí binární
program přes \texttt{cargo run} a~ověřují vznik výstupu v~adresáři \texttt{.build/}.
Tím se simuluje běžné použití kompilátoru z~příkazové řádky.

Všechny testy lze spustit příkazem:
\begin{lstlisting}[language=bash, caption={Spuštění testů}]
cargo test
\end{lstlisting}

\subsection{Výsledky testování}

Testovací sada pokrývá základní syntaktické konstrukce jazyka, práci s~datovým
zásobníkem, překlad uživatelských definic, vybrané řídicí struktury a~optimalizační
průchody nad IR.

V~adresáři \texttt{tests/} jsou testy rozděleny do tematických souborů:
\texttt{lexer\_tests.rs}, \texttt{parser\_tests.rs}, \texttt{analyzer\_tests.rs},
\texttt{ir\_tests.rs}, \texttt{codegen\_tests.rs} a~\texttt{integration\_tests.rs}.
Tyto testy společně ověřují správnost výstupů jednotlivých fází i~to, že
generování kódu odpovídá očekávané sémantice.

Při vývoji byly odhaleny a~opraveny zejména chyby související s~hranami
tokenizace (pozice tokenů přes víceřádkový vstup, escapování řetězců) a~se
správným zacházením s~kontextem REPL (registrace slov a~proměnných mezi vstupy).

\subsection{Porovnání s existujícími implementacemi}

Cílem práce není nahrazovat zavedené implementace jako Gforth, ale popsat
a~ověřit architekturu kompilátoru založenou na explicitním mezikódu,
optimalizačních průchodech a~více backendech.

Ve srovnání s~tradičním \Forth{} je \\Roth{} koncipován především jako
kompilátorový projekt: jazyková sada je omezena na konstrukce potřebné pro
demonstraci frontendu, middle-endu i~backendů. Těžiště práce proto neleží
v~úplné kompatibilitě se standardem \Forth{}, ale v~návrhu překladového řetězce.



%% ===================================================================
%% KAPITOLA 6: ZÁVĚR
%% ===================================================================
%% ===================================================================
%% KAPITOLA 6: ZÁVĚR
%% ===================================================================
\begin{kiconclusions}
	Tato práce popsala návrh a~implementaci kompilátoru zásobníkového jazyka
	inspirovaného jazykem \Forth{}. Výsledná implementace \\Roth{} v~jazyce Rust
	zahrnuje lexikální a~syntaktickou analýzu, sémantickou kontrolu, převod do
	mezikódu, optimalizační průchody a~generování cílového kódu.

	Součástí návrhu je zásobníkový mezikód a~sada optimalizačních průchodů
	(constant folding, inlining, peephole optimalizace, strength reduction
	a~eliminace mrtvého kódu). Implementace podporuje generování kódu do jazyků
	Rust a~C a~obsahuje REPL prostředí založené na průběžné kompilaci a dynamickém
	načítání knihoven.

	Další rozvoj může směřovat k~rozšíření backendů, doplnění dalších instrukcí
	mezikódu, rozšíření standardní knihovny a~k~systematičtějšímu měření výkonu
	na větší sadě referenčních programů.
\end{kiconclusions}

\begin{kiconclusions}[english]
	This thesis described the design and implementation of a stack-based language
	inspired by Forth. The resulting Rust implementation of Roth includes lexing,
	parsing, semantic validation, lowering to an intermediate representation,
	IR-level optimizations, and target code generation.

	The design uses a dedicated stack-based intermediate representation together
	with optimization passes such as constant folding, inlining, peephole
	optimization, strength reduction, and dead code elimination. The implementation
	supports Rust and C backends and includes a REPL environment based on
	incremental compilation and dynamic library loading.

	Future work may focus on extending backend coverage, adding more IR
	instructions, expanding the standard library, and performing more systematic
	performance measurements on a broader benchmark set.
\end{kiconclusions}



%% ===================================================================
%% PŘÍLOHY
%% ===================================================================
\appendix
\section{Obsah elektronických dat}
\label{sec:ObsahData}

\begin{description}
	\item[\texttt{thesis/}] \hfill \\
	      Zdrojové soubory práce v \LaTeX{}u a výsledné PDF (\texttt{kidiplom.pdf}).

	\item[\texttt{src/}] \hfill \\
	      Zdrojové kódy kompilátoru \Roth{} v jazyce Rust (frontend, IR, optimalizace, codegen).

	\item[\texttt{src/repl/}] \hfill \\
	      Implementace REPL režimu s JIT-like kompilací.

	\item[\texttt{roth-runtime/}] \hfill \\
	      Knihovna poskytující běhový kontext (zásobník, paměť) pro REPL.

	\item[\texttt{std/}] \hfill \\
	      Standardní knihovna jazyka ve formě \Roth{} zdrojových souborů.

	\item[\texttt{tests/}] \hfill \\
	      Jednotkové a integrační testy kompilátoru.

	\item[\texttt{test\_source/}] \hfill \\
	      Ukázkové vstupy pro integrační testy (včetně očekávaných výstupů).

	\item[\texttt{README.md}] \hfill \\
	      Stručný popis projektu.

	\item[\texttt{Cargo.toml}] \hfill \\
	      Konfigurační soubor projektu pro Rust.

	\item[\texttt{.build/}] \hfill \\
	      Výstupní adresář generovaného kódu (vytváří se při běhu kompilátoru).
\end{description}

\section{Uživatelská dokumentace}
\label{sec:dokumentace}

Tato příloha popisuje použití kompilátoru \Roth{} z příkazové řádky a základní
práci v režimu REPL.

\subsection{Požadavky}

Pro sestavení projektu je potřeba nainstalovaný Rust toolchain.
Pro spuštění některých režimů jsou dále potřeba:
\begin{itemize}
	\item \texttt{rustc} v \texttt{PATH} (používá se při \texttt{--run} a v režimu REPL),
	\item \texttt{gcc} (používá se při \texttt{--run} s backendem pro C).
\end{itemize}

\subsection{Sestavení a testy}

Projekt lze sestavit a otestovat standardními příkazy:
\begin{lstlisting}[language=bash, caption={Sestavení a testy}]
cargo build
cargo test
\end{lstlisting}

\subsection{Kompilace souboru}

Kompilace souboru probíhá spuštěním binárky \texttt{roth} se jménem vstupního
souboru. Výstupní kód je uložen do adresáře \texttt{.build/}.

\begin{lstlisting}[language=bash, caption={Kompilace souboru}]
cargo run -- path/to/program.fs
\end{lstlisting}

Volba backendu:
\begin{lstlisting}[language=bash, caption={Volba backendu}]
cargo run -- path/to/program.fs --backend rust-ir
cargo run -- path/to/program.fs --backend c-ir
\end{lstlisting}

Volitelně lze nastavit jméno výstupního souboru (stále bude umístěn do
\texttt{.build/}):
\begin{lstlisting}[language=bash, caption={Vlastní název výstupu}]
cargo run -- path/to/program.fs -o out.rs
\end{lstlisting}

\paragraph{Režim \texttt{--run}.}
Při použití \texttt{--run} kompilátor kromě vygenerování kódu automaticky spustí
překlad cílového kódu a výsledný program vykoná. Pro Rust backend se používá
\texttt{rustc -O}, pro C backend \texttt{gcc -O2}.
\begin{lstlisting}[language=bash, caption={Kompilace a spuštění}]
cargo run -- path/to/program.fs --backend rust-ir --run
cargo run -- path/to/program.fs --backend c-ir --run
\end{lstlisting}

\subsection{Ladění výpisů}

Volba \texttt{--debug} řídí množství ladicích informací:
\begin{itemize}
	\item \texttt{0} --- bez výpisů,
	\item \texttt{1} --- vypíše vygenerovaný cílový kód,
	\item \texttt{2} --- vypíše tokeny, AST, IR a statistiky optimalizací,
	\item \texttt{3} --- navíc vypíše barevně zvýrazněný cílový kód (pokud je dostupný highlighter).
\end{itemize}

Volba \texttt{--no-color} vypne barevné zvýraznění.

\subsection{REPL režim}

REPL se spustí bez argumentu vstupního souboru, nebo pomocí \texttt{-i/--interactive}:
\begin{lstlisting}[language=bash, caption={Spuštění REPL}]
cargo run
cargo run -- -i
\end{lstlisting}

V REPL lze psát běžné \Roth{} příkazy (např. \texttt{2 3 + .}) a definovat slova
pomocí \texttt{: NAME ... ;}. REPL navíc obsahuje příkazy začínající dvojtečkou
bez mezery (např. \texttt{:help}).

\begin{lstlisting}[language=bash, caption={Základní REPL příkazy}]
:help      (napoveda)
:stack     (vypis zasobniku)
:words     (seznam uzivatelskych slov)
:vars      (seznam promennych)
:clear     (vypradneni zasobniku)
:reset     (reset stavu)
:debug 2   (nastaveni debug urovne)
:quit      (ukonceni)
\end{lstlisting}

Konstrukce \texttt{INCLUDE} je podporována pouze při kompilaci ze souboru
(je implementována jako preprocesor v \texttt{src/main.rs}); v režimu REPL se
nevyhodnocuje.

\section{Gramatika jazyka}
\label{sec:gramatika}

Tato sekce popisuje implementovanou syntaxi jazyka \Roth{} ve zjednodušené EBNF.
Jazyk je tokenizován na základě bílých znaků a několika speciálních znaků
(\texttt{:}, \texttt{;}, \texttt{(}, \texttt{)}, \texttt{"}).

\subsection{Lexikální pravidla}

\begin{description}
	\item[Whitespace] mezery, tabulátory a konce řádků oddělují tokeny.
	\item[Comment] \texttt{(} \emph{libovolné znaky bez \texttt{)}} \texttt{)}. Komentáře jsou ignorovány parserem.
	\item[Number] dekadický literál \texttt{i32}. Pokud převod selže (např. přetečení), token je brán jako slovo.
	\item[Word] posloupnost znaků do prvního bílého znaku, \texttt{(}, \texttt{)}, \texttt{"}. Slova jsou case-insensitive (lexer je normalizuje na velká písmena).
	\item[String] buď \texttt{"..."} (s podporou escape sekvencí \texttt{\\n}, \texttt{\\t}, \texttt{\\r}, \texttt{\\"}, \texttt{\\\\}), nebo \texttt{S" ..."} (bez escape sekvencí; volitelná mezera po \texttt{S"}).
\end{description}

\subsection{Syntaktická pravidla (EBNF)}

\begin{lstlisting}[language=Forth, caption={EBNF gramatika}]
program        ::= { statement } ;

statement      ::= number
                | string
                | word
                | definition
                | variable_decl
                | comment ;

definition     ::= ':' word { stmt_in_def } ';' ;
stmt_in_def    ::= number | string | word | variable_decl | comment ;
                // vnorene definice nejsou dovoleny

variable_decl  ::= 'VARIABLE' word ;

number         ::= [ '-' ] digit { digit } ;

string         ::= '"' { char | escape } '"'
                | 'S"' [ ' ' ] { char } '"' ;

comment        ::= '(' { char } ')' ;

word           ::= token ;
\end{lstlisting}

Při kompilaci ze souboru je před lexikální analýzou proveden preprocesing
direktivy \texttt{INCLUDE}, která rekurzivně vkládá obsah souborů do zdrojového
textu. Název souboru může být uveden v uvozovkách nebo bez nich a cesta je
vyhodnocena relativně k aktuálnímu souboru (a případně k pracovnímu adresáři).




%% ===================================================================
%% BIBLIOGRAFIE
%% ===================================================================
\nocite{*}
\printbibliography

\end{document}
